\documentclass{beamer}

\input{../helper}

\title{Introdução à Computação \\ operação com strings}
\author{Alexandre Rademaker}
\date{}

\begin{document}

\begin{frame}
 \maketitle
\end{frame}


\begin{frame}[fragile,allowframebreaks]{o tamanho de uma string}

  Podemos calcular o tamanho de uma string:

  \lstinputlisting[style=normalc,caption=src/length0.c]{src/length0.c}

  \framebreak

  E criar uma função para realizar esta tarefa para qualquer string:

  \lstinputlisting[style=normalc,caption=src/length1.c]{src/length1.c}

  \framebreak

  Mas alguém já implementou isso! função \ilcode{strlen} da
  \ilcode{string.h}!

  \lstinputlisting[style=normalc,caption=src/length2.c]{src/length2.c}

  \framebreak

  Veja
  \href{https://man7.org/linux/man-pages/man3/string.3.html}{string
    man page} ou digite `man string' (Linux ou MacOS).

\end{frame}


\begin{frame}[fragile,allowframebreaks]{imprimindo strings}

  Podemos imprimir uma string caractere por caractere. 

  \lstinputlisting[style=normalc,caption=src/string0.c]{src/string0.c}

  \framebreak

  Como tornar o código mais eficiente? Não precisamos recalcular a
  cada interação o tamanho da string!

  Veja que `strlen' agora é usada apenas uma vez na declaração de
  \ilcode{n}.

  \framebreak

  \lstinputlisting[style=normalc,caption=src/string1.c]{src/string1.c}

\end{frame}


\begin{frame}[fragile,allowframebreaks]{transformando em maiúsculas}

  Vamos implementar um programa para transformar caracteres em
  minúsculas para maiúsculas?

  \framebreak

  \lstinputlisting[style=normalc,caption=src/uppercase0.c]{src/uppercase0.c}

  \framebreak

  \includegraphics[width=\textwidth]{ascii-tab.png}

  \framebreak

  Novamente, podemos \emph{reusar} códigos implementados na
  \ilcode{ctype.h}.

  \lstinputlisting[style=ssc,caption=src/uppercase1.c]{src/uppercase1.c}

  \framebreak

  E podemos simplificar.  Na computação, sempre que possível devemos
  \emph{reusar}! Mas existem desvantagens em alguns casos. Quais?

  \lstinputlisting[style=ssc,caption=src/uppercase2.c]{src/uppercase2.c}

    
\end{frame}


\end{document}

%%% Local Variables:
%%% mode: latex
%%% TeX-master: t
%%% End:
